\section{Introduction}

	Reading is often described as "the process of gaining meaning from print" \citep[][p.~34]{raynerHowPsychologicalScience2001}. Getting from a visual word form to its meaning is not a single step: theories of language comprehension generally agree that this process unfolds across multiple stages, moving from sublexical orthographic analysis toward lexical and then semantic representation, though they disagree sharply about how many stages are involved and what mediates between them. English morphology offers a useful case for studying this question. English spelling maps unevenly onto sound, so phonological decoding alone cannot support fluent word recognition \citep{zieglerReadingAcquisitionDevelopmental2005, castlesHowDoesOrthographic2006}; but English morphology is comparatively regular, and words that share a base often show both consistent orthographic patterns and systematic semantic relationships (e.g., \textit{heal--health} share a stem; \textit{teacher}, \textit{baker}, and \textit{driver} share the agentive suffix \textit{-er}). Multi-morphemic words therefore provide a natural test case for asking how form-based and meaning-based information are related across the stages of word recognition.

	Readers vary considerably in how they navigate these stages \citep{perfettiReadingAbilityLexical2007, castlesHowDoesOrthographic2006, stanovichExposurePrintOrthographic1989}, and this variation offers a way to adjudicate between competing accounts of the stages themselves. If a lexical manipulation's effect depends systematically on the reader, that dependency is informative about what kind of representation the manipulation engages. The present study applies this logic to morphological processing. We examine two dissociable dimensions of reading skill, Orthographic Sensitivity (OS) and Language Proficiency (LP), in relation to two lexical variables long argued to index distinct levels of morphological representation: base frequency and morphological family size. By linking these individual differences to the N250 and N400 --- ERP components argued to index sequential stages from sublexical form analysis to semantic integration --- we ask whether the OS/LP dissociation tracks the form-based/meaning-based dissociation in real time. This approach is consistent with the major theoretical accounts of morphological processing: although these accounts disagree about mechanism, they converge on the prediction that base frequency and family size are temporally dissociable and should be differentially sensitive to a reader's orthographic versus semantic profile.

	To evaluate this prediction, four things need to be established before the study's specific aims can be stated. First, what base frequency and family size are each argued to index, and why competing theoretical accounts converge on expecting them to dissociate temporally. Second, how OS and LP emerge as dissociable dimensions of individual variation. Third, why the N250 and N400 are appropriate neural indices of the two processing stages in question. Fourth, what the dissociation between OS and LP predicts specifically for the parsing of novel morphological forms, building on prior evidence that affix-based and stem-based parsing mechanisms may operate concurrently.

	\subsection{Base Frequency and Morphological Family Size}

	Base frequency, the cumulative token frequency of a base morpheme across all its surface forms, and morphological family size, the number of distinct words sharing that base, both facilitate visual word recognition, and their effects are independent when appropriately controlled \citep{fordDerivationalMorphologyBase2010}. The two variables have been argued to index different levels of morphological representation: base frequency reflects the cumulative ease with which a base morpheme's orthographic form is processed across its surface variants, while family size reflects the breadth of meaning-level connections activated when a morphological base is encountered \citep{schreuderHowComplexSimplex1997, taftMorphologicalDecompositionReverse2004, taftRecognitionAffixedWords1979, dejongMorphologicalFamilySize2000, bertramEffectsFamilySize2000}.

	Current accounts of morphological processing disagree considerably about the architecture responsible for this pattern --- whether complex words are obligatorily decomposed into constituent morphemes, accessed via parallel whole-word and decompositional routes, or processed through learned form-to-meaning mappings without discrete morpheme representations at all \citep{schreuderModelingMorphologicalProcessing1995, dejongMorphologicalFamilySize2000, taft_lexical_1975, taftMorphologicalDecompositionReverse2004, rastleMorphologicalDecompositionBased2008, longtinMorphologicalDecompositionEarly2005, morrisSemanticTransparencyMasked2007, chuangDiscriminativeLearningLexicon}. Despite this disagreement, these accounts converge on a shared prediction: base frequency effects should arise earlier in processing than family size effects, reflecting a general precedence of form-level access over semantic integration. We return to how each account explains this precedence, and to what our results imply for adjudicating between them, in the General Discussion; here, the point is simply that the precedence itself is a point of theoretical consensus, which is what motivates using it to test the timing of OS and LP effects.

	\subsection{Orthographic Sensitivity and Language Proficiency as Dissociable Dimensions}

	Traditional models of word recognition typically characterise processing at the level of the skilled reader in general, abstracting over individual differences in favour of a uniform cognitive architecture. Yet readers do not differ along a single dimension of lexical skill. Andrews and colleagues \citep{andrewsIndividualDifferencesSkilled2012, andrewsMorphologicalPrimingStronger2013} showed that individual differences in orthographic and semantic knowledge are dissociable and predict distinct patterns of morphological processing: readers with a semantic profile, defined by relatively higher vocabulary than spelling ability, show robust priming for transparent morphological pairs but little priming for opaque or form-only pairs, particularly at longer latencies, while readers with an orthographic profile, defined by relatively higher spelling than vocabulary, show sustained priming for opaque pairs that is at least as strong as for transparent pairs. \citet{andrewsMorphologicalPrimingStronger2013} interpret this as evidence that readers differ not merely in overall processing efficiency but in the relative weight they assign to form-based versus meaning-based information during lexical access.

	The present study builds directly on this finding. If base frequency and family size index form-based and meaning-based aspects of lexical representation respectively, as established above, then readers who differ in their balance of orthographic and semantic processing should differ predictably in their sensitivity to each variable. To capture this variation continuously rather than as a categorical split, we derive two orthogonal dimensions, Language Proficiency and Orthographic Sensitivity, from a principal components analysis of vocabulary knowledge, spelling ability, and print exposure, following \citet{andrewsMorphologicalPrimingStronger2013}.

	\subsection{The N250 and N400 as Indices of Sequential Processing Stages}

	To track whether base frequency and family size effects emerge at distinct temporal stages, the present study uses two ERP components as dependent measures: the N250, an index of sublexical form processing, and the N400, an index of lexical-semantic integration. These components have been argued to reflect sequential stages in the cascade from orthographic decoding to meaning retrieval \citep{graingerWatchingWordGo2009}.

	The N250 is a negative-going deflection peaking at approximately 250~ms post-stimulus onset, proposed to reflect the stage at which sublexical orthographic representations are mapped onto higher-level form representations prior to full lexical access \citep{holcombTimeCourseVisual2006}. Evidence for this interpretation comes from its graded sensitivity to orthographic overlap in masked priming: amplitude is most attenuated for full identity repetitions, intermediate for primes sharing all but one letter, and largest for unrelated pairs \citep{holcombTimeCourseVisual2006, holcomb_exploring_2007}. The N250 is also sensitive to sublexical phonological overlap, with separable orthographic and phonological subcomponents \citep{graingerTimeCourseOrthographic2006}, consistent with its role as an early, form-based stage of processing.

	The N400 is a negative-going component peaking around 400~ms with a broad centro-parietal distribution \citep{kutasReadingSenselessSentences1980}. Its amplitude varies with semantic priming, cloze probability, and semantic richness, while remaining insensitive to physical and syntactic manipulations that leave meaning intact \citep{kutasThirtyYearsCounting2011}. In the present design, larger amplitudes are expected when a word's meaning is harder to integrate with the concurrently activated semantic information carried by its morphological family.

	The N250 and N400 are dissociable not only by topography --- frontal versus posterior \citep{holcombTimeCourseVisual2006} --- but by a double dissociation across priming types: lexical status modulates the N400 but not the N250, since N400 priming is restricted to real words and does not extend to pseudowords \citep{kiyonagaMaskedCrossmodalRepetition2007}, while the reverse holds for semantic relatedness, which modulates the N400 but produces no N250 effect \citep{graingerWatchingWordGo2009}. This pattern confirms that the two components index qualitatively distinct processing stages, ruling out the interpretation of the N250 as merely an early onset of the N400.

	\subsection{Affix-Led and Stem-Led Parsing}

	A recurring question in this literature is whether morphological decomposition is driven primarily by affix recognition or by stem activation. Affix-stripping accounts \citep{taft_lexical_1975, rastleMorphologicalDecompositionBased2008} propose that the affix is identified first and triggers decomposition, after which the stem is looked up in the lexicon. However, \citet{beyersmannEmbeddedStemPriming2016} found that high-proficiency readers showed embedded stem priming effects independent of affix status --- priming occurred even when the stem was followed by a non-affixal ending --- suggesting that stem activation in skilled readers need not be contingent on affix recognition, and that stem-based and affix-based mechanisms may operate concurrently.

	The present study tests this possibility directly. If Orthographic Sensitivity indexes the precision of sublexical orthographic representations, readers higher in Orthographic Sensitivity should show greater N250 sensitivity to the presence of a real versus pseudo-affix, consistent with affix-led parsing. If Language Proficiency instead indexes the depth and breadth of lexical-semantic knowledge in which stems are embedded, readers higher in Language Proficiency should show a distinct signature of stem-led parsing at the same stage. Because Beyersmann et al.'s proficiency measure was not decomposed into orthographic and lexical-semantic components, it is unclear from their data alone which mechanism their finding reflects. The present design, with Orthographic Sensitivity and Language Proficiency measured as separate dimensions, is positioned to address this question directly.

	\subsection{Aims and Predictions}

	The present study has two complementary aims. The first is to establish whether base frequency and morphological family size exert their effects at distinct temporal stages of visual word recognition. The second, and primary, aim is to examine whether Orthographic Sensitivity and Language Proficiency modulate these effects selectively, at the processing stages where each variable is predicted to operate.

	The core prediction is a pattern of component-specific modulation: Orthographic Sensitivity should modulate base frequency effects at the N250, while Language Proficiency should modulate family size effects at the N400. If this pattern is observed, it would constitute direct electrophysiological evidence that the distinction between form-based and meaning-based morphological processing is a real axis of variation in the reading system, not merely a property of competing computational models. If instead both variables are modulated by the same individual difference dimension at the same stage, this would suggest that the two effects are less dissociable than current multi-stage accounts propose.

	The first analysis examines this prediction in real word recognition. Because word stimuli were not matched across morphological variables, base frequency, family size, and morphological complexity are treated as continuous predictors in a regression framework for reaction times; for ERP data, items were median-split into high and low groups separately for base frequency and family size. The two measures therefore carry somewhat different inferential weight: the RT regression speaks to the independent continuous contributions of each variable, while the ERP analyses speak to the direction and timing of effects at the extremes of each variable's distribution.

	The second analysis examines whether these modulations generalise to novel morphological structure, using morphologically complex nonwords. Morphologically structured nonwords (e.g., \textit{gasful}: real stem + real affix) were compared against matched controls (e.g., \textit{gasfil}: real stem + pseudoaffix), holding stem identity constant while manipulating suffix legitimacy. Because these items have no stored whole-word representations, they place maximal pressure on the system's generalisation mechanism and allow stem-level variables to be examined independently of whole-word familiarity. Orthographic Sensitivity is predicted to modulate early N250 responses to the suffix contrast, consistent with affix-led parsing: readers with more precise representations of real suffix forms should show greater N250 sensitivity to the distinction between real and pseudo-suffixes. Language Proficiency, building on the stem-led parsing possibility developed above, is predicted to modulate N250 responses through a distinct mechanism, and to predict later N400 sensitivity to whether a nonword's stem belongs to a large or productive morphological family, since it is this family membership that determines the richness of the semantic activation the form can support.